% !TeX program = xelatex
% Scientific code: 4c1ccd79; oral-readout origin: 8019f29e.
\documentclass[12pt]{article}
\usepackage[a4paper,margin=2.5cm]{geometry}
\usepackage{amsmath,amssymb,bm}
\usepackage{booktabs,tabularx,array,longtable}
\usepackage{fontspec}
\usepackage{needspace}
\usepackage[unicode,hidelinks]{hyperref}
\defaultfontfeatures{Ligatures=TeX}
\setmainfont[AutoFakeSlant=0.2]{Noto Serif CJK SC}
\setsansfont{Noto Sans CJK SC}
\setmonofont{Noto Sans Mono CJK SC}
\XeTeXlinebreaklocale "zh"
\XeTeXlinebreakskip=0pt plus 1pt minus 0.1pt
\setlength{\parindent}{2em}
\setlength{\parskip}{0.25em}
\setlength{\emergencystretch}{3em}
\setlength{\LTpre}{0.6em}
\setlength{\LTpost}{0.6em}
\renewcommand{\arraystretch}{1.2}
\allowdisplaybreaks
\setcounter{tocdepth}{1}
\renewcommand{\contentsname}{目录}
\makeatletter
\renewcommand{\l@section}{\@dottedtocline{1}{0em}{1.8em}}
\makeatother
\newcommand{\ind}{\mathbb I}
\newcommand{\logit}{\operatorname{logit}}
\newcommand{\sigmoid}{\operatorname{sigmoid}}
\newcommand{\clip}{\operatorname{clip}}
\newcommand{\R}{\mathcal R}
\newcolumntype{L}[1]{>{\raggedright\arraybackslash}p{#1}}
\title{Model 0923：有限候选、反馈学习与口述表达\\[0.4em]
\large 以认知模块为线索的完整模型说明}
\author{}
\date{2026年10月10日}
\begin{document}
\maketitle

本文描述已经接入代码的 \texttt{0923-M1-v1}。目标是用更容易解释的机制，保持已有模型的行为预测能力。目前已完成实现测试和三条件短序列检查，还没有完成 M1 的被试参数拟合、参数恢复和表现比较。下面说明模型怎样工作，以及各项设计意味着什么，不把设计假设写成已经证实的认知结论。

正文按照“看到刺激、想到规则、做出选择、说出判断、收到反馈、准备下一题”的顺序组织。第5节和第6节分别说明搜索控制与候选管理，第9节完整说明 oral 模块。被试参数在同一被试内保持不变，状态随试次递推。

\tableofcontents
\clearpage

\section{模型概述：几个认知模块怎样一起工作}
\label{sec:overview}

模型假设，学习者每次只考虑少量规则。他可能同时保留“看脖子长短”和“比较脖子与尾巴”这两种解释，却在当前题主要使用其中一条。收到反馈后，他会重新评价这些解释，也可能想到一条新规则，或不再考虑某条弱规则。候选数量可以变化，容量只规定最多能同时保留多少条。

一条规则有两个不同的状态。规则信念表示“在当前候选中，更相信哪条”；规则精度表示“如果按这条规则判断，类别预测有多明确”。执行规则则是本题实际控制作答的那一条。高信念通常提高执行机会，但执行还受上一题使用习惯影响；高精度使预测更明确，却不保证规则正确。

口述来自同一次内部判断。学习者先作答，再描述本题判断，之后才收到外部反馈。他可以少说一个条件，也可以顺便描述眼前刺激的另一个部位。因此，模型既不要求口述等于完整规则，也不把每个被提到的部位都当作分类依据。

各模块的分工如下。这里按认知问题划分，不要求每一行都对应独立的 Python 类。

{\small
\begin{longtable}{L{2.6cm}L{5.6cm}L{6.0cm}}
\toprule
模块 & 要回答的问题 & 主要输出\\
\midrule
\endhead
知觉模块 & 被试看到了怎样的刺激？ & 内部刺激 $\widetilde{\bm x}_t$。\\
规则预测 & 某条规则会怎样分类这幅刺激？ & 类别概率 $p_{t,h}(c)$。\\
证据记忆 & 反馈怎样改变对候选的信任？ & 反馈后信念 $\pi_t^+$；三值任务还更新配对知识。\\
搜索控制 & 下一题有多大机会想到新规则？偏局部还是全局？ & 失败状态 $F_{t+1}$、搜索概率 $E_{t+1}$、全局比例 $g_{t+1}$。\\
候选管理 & 哪些规则继续保留？新人怎样入场并获得支持？ & 集合 $A_{t+1}$ 和信念 $\pi_{t+1}^-$。\\
精度学习 & 反馈使规则预测更明确还是更模糊？ & 精度 $\beta_{t+1}(h)$。\\
执行与选择 & 本题用哪条规则，选哪个类别？ & 执行规则 $e_t$ 与选择 $y_t$。\\
口述模块 & 怎样把本题判断说出来？ & 给定选择的报告分布 $P(o_t\mid y_t)$。\\
\bottomrule
\end{longtable}
}

粒子滤波用于推断这些隐藏过程。每个粒子是一条可能的内部学习路径，多个粒子共同表达研究者的不确定性。粒子滤波不是额外的认知模块，也不表示被试在脑中同时运行这些粒子。

\section{任务输入、规则目录与潜在状态}
\label{sec:states}

\subsection{每题输入什么}

对被试 $s$ 的第 $t$ 题，物理刺激、实际选择、外部反馈和可用口述分别记为
\begin{equation}
\bm x_{s,t}\in[0,1]^4,\qquad
y_{s,t}\in\{1,\ldots,C\},\qquad r_{s,t},\qquad o_{s,t}.
\end{equation}
四维刺激按该被试任务的实际特征顺序排列。下文省略被试下标。Condition 1 为二分类二值反馈，$C=2$，$r_t\in\{0,1\}$；condition 2 为四分类二值反馈，$C=4$，$r_t\in\{0,1\}$；condition 3 为四分类三值反馈，$r_t\in\{0,0.5,1\}$。

实际选择用于计算行为分数和修正粒子权重，选择与反馈随后用于被试模型的学习。口述已有显式生成接口，但还没有进入批量行为拟合的粒子权重。反应时不进入本版模型。

真实目标规则不进入搜索或学习控制器；真实类别用于环境评分和结果评价。Condition 3 的正确按键配对也不作为被试的初始知识。

\subsection{固定规则目录}

目录记为 $\mathcal H=\{1,\ldots,J\}$。Condition 1 有29条规则：

{\small
\begin{longtable}{L{3.7cm}L{8.6cm}L{1.3cm}}
\toprule
规则形式 & 判断例子或边界 & 数量\\
\midrule
\endhead
单特征阈值 & $x_i\leq0.5$，按某个部位的长短分类。 & 4\\
两个特征的顺序 & $x_i\leq x_j$，比较脖子和尾巴。 & 6\\
两个特征的和 & $x_i+x_j\leq1$。 & 6\\
两组特征和的顺序 & $x_i+x_j\leq x_k+x_l$。 & 3\\
两个特征是否接近 & $|x_i-x_j|\leq0.10$。 & 6\\
单特征是否在中间 & $|x_i-0.5|\leq0.10$。 & 4\\
\bottomrule
\end{longtable}
}

Condition 2 和3使用128条四分类规则，包括两个判断的交叉组合、条件分支、最大／最小特征，以及允许一个特征出现在两个判断中的12条组合规则。完整分组见附录\ref{app:catalogue}。三条件使用各自目录，不能把29条二分类规则直接当作四分类规则。

每条规则的类别区域及标签固定。当前模型没有另加区域与类别标签的关联学习，也不为提高相似度重新排列标签。目录规定模型允许考虑哪些解释，不表示被试一定知道或用过所有这些规则。

基础先验在相应目录内均匀：
\begin{equation}
p_0(h)=1/J,\qquad h\in\mathcal H.
\label{eq:base-prior}
\end{equation}
它用于初始化、全局提出和功能提出的基础权重，不按后来观察到的学习类型提前偏好规则。

\subsection{容量上限与初始候选}

容量参数为 $M_s^{\max}$，实际集合为 $A_t$，数量为 $N_t$：
\begin{equation}
A_t\subset\mathcal H,\qquad N_t=|A_t|,\qquad
1\leq N_t\leq M_s^{\max}.
\label{eq:capacity}
\end{equation}
集合不需要填满。第一题从 $p_0$ 抽取一条规则，令 $A_1=\{h_1\}$，并用它执行。当前要求整数 $2\leq M_s^{\max}<J$，以允许执行规则和备选共存，并保留非活跃候选；实际数量仍可以为1。

只保留一条候选时，其候选内信念为1，表示“目前只有这个解释在比较”，不表示“已经确定它正确”。反馈是否支持它仍能体现在精度和失败状态中。初始一条候选及容量取值约束是模型约定，尚未由行为数据证明。

\subsection{需要保存哪些状态}

每条内部路径在作答前保存
\begin{equation}
z_t=\bigl(\widetilde{\bm x}_t,A_t,\pi_t^-,F_t,
\bm\beta_t,e_t,\Pi_t^-\ \text{（仅condition 3）}\bigr).
\end{equation}
$\pi_t^-$ 在活跃集合内归一化，非活跃概率为0；$e_t\in A_t$；$\Pi_t^-$ 为规则与配对的联合信念。口述是当前判断的观测，不另存会反过来控制学习的“口述规则”。

本文统一使用以下术语和符号：

{\small
\begin{longtable}{L{2.0cm}L{12.0cm}}
\toprule
符号 & 含义\\
\midrule
\endhead
$\pi_t^-,\pi_t^+$ & 作答前、反馈后的相对规则信念。\\
$\beta_t(h)$ & 规则把类别距离转成概率时使用的精度。\\
$e_t$ & 当前实际执行规则。\\
$F_t$ & 进入当前题的失败状态，不表示规则证据。\\
$E_t,g_t$ & 搜索启动概率、提出时的全局来源比例。\\
$\gamma_s,\delta_s$ & 分别控制旧规则证据、失败状态的保持。\\
$\omega_s$ & 局部搜索中选择结构编辑的倾向。\\
$q_t(k\mid h)$ & Condition 3 的条件配对信念，不与 $\omega_s$ 混用。\\
\bottomrule
\end{longtable}
}

\section{知觉模块：从刺激到单条规则的预测}
\label{sec:perception}

\subsection{形成带有误差的内部刺激}

同一物理刺激在不同试次可以引起略有差异的知觉：
\begin{equation}
\widetilde x_{t,j}=\clip(x_{t,j}+\varepsilon_{s,t,j},0,1).
\label{eq:perception}
\end{equation}
误差参数来自学习任务前的独立知觉测量，并按实际特征顺序排列。当前个体化配置通常使用测量得到的逐维正态误差；共享模块也允许按测量定义使用均匀误差。具体分布与参数以运行配置为准。

同一内部刺激用于规则预测与口述补充描述，不为口述重新感知一次。粒子使用独立且可复现的知觉随机流，汇总不可直接观察的知觉差异。

\subsection{单条规则怎样给出类别概率}

规则 $h$ 对应类别区域 $R_{h,1},\ldots,R_{h,C}$。先算内部刺激到区域的最短欧氏距离：
\begin{equation}
d_{t,h,c}=\inf_{\bm z\in R_{h,c}}
\|\widetilde{\bm x}_t-\bm z\|_2.
\label{eq:distance}
\end{equation}
区域内部距离为0；一个类别有多个分离部分时取最近部分。随后
\begin{equation}
p_{t,h}(c)=
\frac{\exp[-\beta_t(h)d_{t,h,c}]}
{\sum_{a=1}^{C}\exp[-\beta_t(h)d_{t,h,a}]}.
\label{eq:emission}
\end{equation}
精度越高，较近类别越占优势；精度越低，预测越接近均匀。知觉误差改变输入，精度改变对距离的反应，作用位置不同。

这一层回答“如果按规则 $h$ 判断，会怎样分类”，还不决定是否相信或执行它。用于行为的精度也用于评价反馈，不另设第二套证据精度。

\section{证据记忆模块：反馈怎样改变规则信念}
\label{sec:memory}

\subsection{规则是否能解释当前反馈}

二值反馈的概率为
\begin{equation}
\ell_t(h)=r_t p_{t,h}(y_t)
+(1-r_t)[1-p_{t,h}(y_t)].
\label{eq:binary-feedback}
\end{equation}
答对支持能预测已选类别的规则；答错支持倾向其他类别的规则。四分类错误只说明所选答案不对，因此概率仍是 $1-p_{t,h}(y_t)$，不能指定某个正确答案。

把 $\ell_t(h)$ 在活跃集合内归一化为 $L_t(h)$ 不改变似然比，在归一化的信念更新中等价。Condition 3 的0分和半分含义不同，使用第\ref{sec:pairing}节的联合概率。

\subsection{保留旧证据，再加入新证据}

反馈后信念为
\begin{equation}
\pi_t^+(h)=
\frac{\ind(h\in A_t)[\pi_t^-(h)]^{\gamma_s}\ell_t(h)}
{\sum_{v\in A_t}[\pi_t^-(v)]^{\gamma_s}\ell_t(v)},
\qquad 0\leq\gamma_s\leq1.
\label{eq:memory}
\end{equation}
对数形式是“保留一部分旧对数信念，再加本题对数证据”。$\gamma_s$ 越大，旧证据保留越久；越小，旧候选差异被压缩得越快。$\gamma_s=1$ 是无额外遗忘的顺序更新；$\gamma_s=0$ 在活跃集合内放弃旧差异，只凭当前反馈比较。

M1 只使用这一条衰减记忆，永久记忆权重为0、反馈增益为1。候选管理先设置新集合的初始支持，此后反馈才加入，因此新人继承的支持与它亲自得到的反馈是两件事。

所有活跃候选都接受同一结果的评价，包括未执行的备选，表示“如果按另一条候选判断，能不能解释这个结果”。非活跃规则不在后台继续学习全部反馈。

\section{搜索控制模块：何时搜索，向哪里搜索}
\label{sec:search}

\subsection{用一条失败状态概括近期经历}

本题失败程度为 $f_t=1-r_t$。二值反馈对应0或1；三值对应0、0.5或1。半分表示部分未达到目标，描述达成程度，不是反馈信息量。

下一题失败状态为
\begin{equation}
F_{t+1}=\delta_sF_t+(1-\delta_s)f_t,\qquad
F_1=0,\qquad 0\leq\delta_s<1.
\label{eq:failure}
\end{equation}
最新反馈立即进入，随后衰减。较小的 $\delta_s$ 使状态对最新反馈变化较大、恢复较快；较大的值使变化较慢、影响更久。例如 $\delta_s=0.6$ 时，从0开始连续两次失败产生0.4、0.64，再成功一次降到0.384。

$\delta_s$ 是被试参数，0.6只是示例。它与证据记忆 $\gamma_s$ 分开：失败后长时间愿意探索，不等于更能记住哪条规则曾得到支持。完整试次顺序连续递推，session 边界不自动清零。

\subsection{失败状态怎样提高搜索机会}

启动概率为
\begin{equation}
E_{t+1}=\sigmoid[\logit(E_{0,s})+a_sF_{t+1}],
\label{eq:search-probability}
\end{equation}
其中 $\sigmoid(v)=1/(1+e^{-v})$，$\logit(p)=\log[p/(1-p)]$。$E_{0,s}\in(0,1)$ 是无失败影响时的基础搜索概率；$a_s\geq0$ 是失败推动搜索的增益。

连续成功使搜索回到基础水平，反复失败提高搜索机会。基础搜索允许已有办法奏效时仍偶尔想到别的解释；$a_s=0$ 时失败不影响启动。

从第二题起，每题只抽一次
\begin{equation}
Z_{t+1}\sim\operatorname{Bernoulli}(E_{t+1}).
\end{equation}
第一题只是初始化。$Z_{t+1}=1$ 表示尝试提出一个新人；新人可以被拒绝，因此搜索不等于入场，也不等于执行切换。

\subsection{失败状态怎样扩大范围}

有搜索时，全局来源比例为
\begin{equation}
g_{t+1}=g_{0,s}+(1-g_{0,s})c_{G,s}F_{t+1},
\qquad g_{0,s},c_{G,s}\in[0,1].
\label{eq:global-range}
\end{equation}
$g_{0,s}$ 是基础全局比例，$c_{G,s}$ 控制失败使提出更偏整个目录的程度。失败为0时 $g=g_0$，为1时 $g=g_0+(1-g_0)c_G$，始终为合法概率。

启动和范围读取同一失败状态，却输出不同控制：$E$ 决定有没有尝试，$g$ 决定局部／全局来源。$g$ 不直接重置信念、调整精度或抹去配对知识。

\subsection{本模块需要估计什么}

本模块有5个被试参数：$E_{0,s},a_s,\delta_s,g_{0,s},c_{G,s}$。$F_t,E_t,g_t$ 是状态或计算结果，不逐题自由拟合。

\section{候选管理模块：保留、想到与接纳规则}
\label{sec:management}

模块读取反馈后的 $A_t,\pi_t^+,e_t$，产生下一题集合与信念。分清旧规则退出、提出新人和新人入场三个步骤：退出不要求搜索，搜索也不保证入场。

\subsection{低支持备选可以退出，不强制补满}

对每条非执行候选独立抽取
\begin{equation}
P(h\text{退出}\mid A_t,\pi_t^+,e_t)
=\rho_s[1-\pi_t^+(h)],\qquad h\neq e_t,\quad\rho_s\in[0,1].
\label{eq:exit}
\end{equation}
执行规则暂受保护，使集合不为空。弱备选更容易退出，$\rho_s$ 决定速度。记忆改变现有候选的支持，退出改变当前可访问的候选，两者分别解释。

令退出集合为 $D_t$，剩余集合为 $R_t=A_t\setminus D_t$。出现空位不自动补新人，剩余支持为
\begin{equation}
\Gamma_t(h)=
\frac{\pi_t^+(h)}{\sum_{v\in R_t}\pi_t^+(v)},\qquad h\in R_t.
\label{eq:remaining-belief}
\end{equation}
无搜索时提交 $R_t$ 和 $\Gamma_t$。极端情况下，剩余支持总和恰为0，代码显式回退到 $p_0$ 在剩余集合上的条件分布并记录此事，只处理未定义的条件化。

新人来自 $I_t=\mathcal H\setminus A_t$，刚退出的同一规则不在本事件立即返回，但以后可以再次想到。退出者不保存可直接恢复的精度，也不学习缺席期间的反馈。

\subsection{局部搜索的两种理由：结构容易改，或分类结果相近}

从“看脖子长短”改为“看尾巴长短”，只更换判断特征，却可能改变一半刺激的分类。所以“容易改出来”和“分类结果相近”分别建模。

\subsubsection{结构提出：一次合法判断编辑}

规则用规范判断程序表示，包含特征、比较原语、参照、分支和固定类别输出。单步编辑可以更换一个特征、替换比较原语及标准参照、修改一个参照值，或重写条件树的一条下级判断。结果仍须是目录中完整合法的规则，其余判断和输出布局按操作定义保留。

设 $\mathcal M(h)$ 为规范单步编辑，$c(m)>0$ 为成本：
\begin{equation}
d_{\rm edit}(h,u)=
\begin{cases}
\min_{m\in\mathcal M(h):\,m(h)=u}c(m),&\text{有合法编辑},\\
\infty,&\text{否则}.
\end{cases}
\end{equation}
等价写法先规范化，不按编号或字符串长度计算。共同成本为特征更换1、比较原语替换1.5、参照修改1、下级分支重写2。它们是初版假设，不是心理测量，也不按被试分别拟合四项成本。

可用邻居为 $\mathcal N_t(h)=\{u\in I_t:d_{\rm edit}(h,u)<\infty\}$。非空时
\begin{equation}
K_{{\rm edit},t}(u\mid h)=
\frac{\ind[u\in\mathcal N_t(h)]e^{-d_{\rm edit}(h,u)}}
{\sum_{v\in\mathcal N_t(h)}e^{-d_{\rm edit}(h,v)}}.
\label{eq:edit-proposal}
\end{equation}
低成本编辑更容易提出；目录不存在的完整结果不能凭空生成。

\subsubsection{功能提出：比较实际分类结果}

在预定刺激参考分布 $\mu_X$ 下记录
\begin{equation}
C_{ab}^{h,u}=\Pr_{\bm x\sim\mu_X}[c_h(\bm x)=a,\ c_u(\bm x)=b],
\end{equation}
\begin{equation}
S(h,u)=\sum_{a=1}^{C}C_{aa}^{h,u},\qquad
d_{\rm func}(h,u)=1-S(h,u).
\label{eq:functional-distance}
\end{equation}
$c_h$ 是固定标签的确定分类函数，$S$ 是相同标签比例。不用新人的未学习精度来定义内容相似，避免把两个模糊输出误当成相同规则。

当前 $\mu_X$ 为 $[0,1]^4$ 上的均匀分布。用100,000个共同样本、固定种子0计算完整张量，再由同一张量得到一致率。这是积分近似，不表示被试经历了十万题。矩阵保留差异来自哪些类别的信息，当前提出与支持分配使用一致率。

一致率受类别大小影响。两条独立中心带规则各把20\%空间判为带内，一致率为 $0.2^2+0.8^2=0.68$，其中0.64来自都判为带外；这不表示判断思路很近。结构与功能共同保留，以免一个分数承担所有含义。

功能提出核为
\begin{equation}
K_{{\rm func},t}(u\mid h)=
\frac{\ind(u\in I_t)p_0(u)e^{-d_{\rm func}(h,u)/\tau_P}}
{\sum_{v\in I_t}p_0(v)e^{-d_{\rm func}(h,v)/\tau_P}},
\qquad \tau_P=0.10.
\label{eq:functional-proposal}
\end{equation}
固定的 $\tau_P$ 越小，提出越偏分类结果相近的规则。当前没有个人注意权重、经验加权距离或动态相似度。

\subsubsection{混合局部方式，接上全局来源}

局部搜索先按 $\Gamma_t(h)$ 选起点，再按 $\omega_s$ 的比例选结构提出：
\begin{equation}
K_{L,t}(u\mid h)=
\omega_sK_{{\rm edit},t}(u\mid h)
+(1-\omega_s)K_{{\rm func},t}(u\mid h).
\end{equation}
没有可用结构邻居时，该起点的结构核回退为功能核并记录回退，仍属于局部来源。随后
\begin{align}
Q_{L,t}(u)&=\sum_{h\in R_t}\Gamma_t(h)K_{L,t}(u\mid h),\\
Q_{G,t}(u)&=\frac{\ind(u\in I_t)p_0(u)}{\sum_{v\in I_t}p_0(v)},\\
Q_t(u)&=(1-g_{t+1})Q_{L,t}(u)+g_{t+1}Q_{G,t}(u).
\label{eq:proposal-mixture}
\end{align}
有搜索时从 $Q_t$ 提出一个新人。$\omega_s$ 控制局部内部方式，$g_{t+1}$ 控制局部／全局比例。提出来源决定怎样想到它，不额外奖励同一候选的支持。

\subsection{支持继承由新人和旧集合的关系决定}

容易想到不等于相信。结构用于提出，功能关系用于支持继承：
\begin{equation}
B_s(h,u)=e^{-d_{\rm func}(h,u)/\tau_{T,s}},\qquad\tau_{T,s}>0.
\label{eq:affinity}
\end{equation}
$\tau_{T,s}$ 是被试支持泛化尺度：小值主要迁移给功能很近的新人，大值允许更远新人继承支持，与固定提出尺度 $\tau_P$ 分别定义。

每个旧来源在自身和新人之间分配：
\begin{equation}
w(h)=\frac{\Gamma_t(h)}{1+B_s(h,u)},\qquad
w(u)=\sum_{h\in R_t}\frac{\Gamma_t(h)B_s(h,u)}{1+B_s(h,u)}.
\label{eq:temporary-weights}
\end{equation}
每个来源分出的总量等于原量，整体仍为1。接近高支持旧规则的新人通常得到更多；只接近弱规则或远离整个集合时较少，不添加固定新人预算。

例如剩余支持为 $(0.8,0.2)$，亲和度为 $(0.5,0.1)$，旧规则暂为0.533和0.182，新人约0.285。亲和度若为 $(0.02,0.02)$，新人约0.020。差异来自关系和信念，不按后来的学习类型设值。

因 $B\leq1$，新人临时质量不超过 $1/2$，需检查能否容纳快速发现。继承是一项泛化假设，不表示新人已经得到未来反馈。

\subsection{有空位加入，满员时竞争}

若 $|R_t|<M_s^{\max}$，新人加入，式\eqref{eq:temporary-weights}成为下一题信念。满员时在非执行旧候选和新人中淘汰最低临时权重者，并列最低则等权抽取；上一执行规则受保护。

新人被淘汰时，搜索已尝试但未入场，旧候选恢复 $\Gamma_t$。旧规则被淘汰时，在最终集合 $A'=A_{t+1}$ 上从原始 $\Gamma_t$ 重新计算支持，不在临时权重上再次扣除。临时竞争可比较上限外的一条提议，实际提交数量仍不超过上限。

对旧来源 $h\in R_t$ 和最终接收者 $v\in A'$，
\begin{equation}
W(v\mid h,A')=
\begin{cases}
1,&h\in A',\ v=h,\\
B_s(h,u),&h\in A',\ u\in A',\ v=u,\\
B_s(h,v),&h\notin A',\ v\in A',\\
0,&\text{其他},
\end{cases}
\end{equation}
\begin{equation}
T(v\mid h,A')=\frac{W(v\mid h,A')}{\sum_{z\in A'}W(z\mid h,A')},
\qquad
\pi_{t+1}^-(v)=\sum_{h\in R_t}\Gamma_t(h)T(v\mid h,A').
\label{eq:transport}
\end{equation}
幸存来源只在自身和新人间分配，容量竞争淘汰者按亲和度分给最终接收者。每个来源质量守恒、最终归一化；无淘汰时退化为临时公式，拒绝新人时为恒等迁移。自主退出则按式\eqref{eq:remaining-belief}条件化，不声称保留退出者的绝对质量。

计算使用稳定归一化，不能为修复指数下溢而悄悄添加固定新人质量。

\subsection{精度和配对怎样衔接}

幸存者保留精度，新人或重新入场者恢复 $\beta_{0,s}$，之后的实际反馈才更新它。Condition 3 先保存迁移前整体配对，幸存者保留自己的条件配对，新人继承整体配对知识；公式见第\ref{sec:pairing-transfer}节。规则质量守恒不等于配对边际一定不变。

\subsection{本模块需要估计什么}

候选管理及衔接的执行选择有5个被试参数：$M_s^{\max},\rho_s,\omega_s,\tau_{T,s},\lambda_s$，执行惯性见第\ref{sec:execution}节。初始一条、单新人、执行保护和最低权重竞争是固定结构约定，不因没有增加参数就成为事实。

\section{精度学习模块：预测怎样变得明确}
\label{sec:precision}

\subsection{实际结果是否比机会水平更符合规则}

精度更新使用反馈前预测。二值任务令 $b_t(h)=\ell_t(h)$，机会概率为
\begin{equation}
b_{\rm chance}(r_t)=\frac{r_t}{C}+(1-r_t)(1-1/C).
\end{equation}
二分类两种反馈均为0.5；四分类正确为0.25、错误为0.75。支持量为
\begin{equation}
\kappa_t(h)=
\begin{cases}
\dfrac{b_t(h)-b_{\rm chance}}{1-b_{\rm chance}},&b_t(h)\geq b_{\rm chance},\\[6pt]
\dfrac{b_t(h)-b_{\rm chance}}{b_{\rm chance}},&b_t(h)<b_{\rm chance}.
\end{cases}
\label{eq:binary-precision-support}
\end{equation}
二分类等价于 $2[b_t(h)-0.5]$。四分类按反馈事件居中，使均匀预测在正确与错误时都中性，避免错误较常见就导致精度漂移。使用绝对反馈概率，不使用在规则间归一化的似然。

\subsection{支持时增加，反驳时降低}

所有活跃规则使用
\begin{equation}
\beta_{t+1}(h)=
\begin{cases}
\min\{\beta_t(h)+\eta_{+,s}\kappa_t(h)[\beta_{\max}-\beta_t(h)],
\beta_{\max}\},&\kappa_t(h)\geq0,\\[4pt]
\max\{\beta_t(h)-\eta_{-,s}\beta_t(h)\min[1,-\kappa_t(h)],
\beta_{\min}\},&\kappa_t(h)<0.
\end{cases}
\label{eq:precision-update}
\end{equation}
共同边界为 $\beta_{\min}=0.1,\beta_{\max}=25$。$\beta_{0,s}$ 是初值，$\eta_{+,s},\eta_{-,s}$ 是增长和下降速率。支持时增长取决于到上限的余量，反驳时按当前值的一定比例下降。

未执行的活跃候选也更新。新精度从下一题影响预测，不反过来重算本题选择。Condition 3 使用第\ref{sec:pairing-precision}节的三值支持量进入同一更新结构。

精度提高表示输出明确，不保证正确。相同信念可有不同精度，相同精度也可有不同信念。两者共用一次反馈，改变不同状态。

\section{执行与选择模块：本题按哪条规则作答}
\label{sec:execution}

\subsection{信念提供基础权重，上一执行有惯性}

每题完成候选管理后都选执行规则，不要求先搜索：
\begin{equation}
\psi_{t+1}(h)=
\frac{\pi_{t+1}^-(h)e^{\lambda_s\ind(h=e_t)}}
{\sum_{v\in A_{t+1}}\pi_{t+1}^-(v)e^{\lambda_s\ind(v=e_t)}},
\qquad e_{t+1}\sim\psi_{t+1}.
\label{eq:execution}
\end{equation}
$\lambda_s\geq0$ 给上一执行规则乘以 $e^{\lambda_s}$，不是永远固定它。两候选信念均为0.5，$\lambda_s=\log3$ 时继续使用上一规则的概率为0.75；备选得到更强支持时仍可切换。

$\lambda_s=0$ 表示执行读出无额外惯性，候选管理中的执行保护仍存在。“想到新人”“新人入场”“开始执行”可以发生在不同题；没有新人也可改用旧备选，新人入场也可暂不执行。

\subsection{由实际执行规则生成选择}

已抽取执行规则的路径中，
\begin{equation}
P(y_t=c\mid z_t)=p_{t,e_t}(c).
\label{eq:path-choice}
\end{equation}
抽取执行前，把它积分掉则为
\begin{equation}
P(y_t=c\mid\text{抽取执行前状态})
=\sum_{h\in A_t}\psi_t(h)p_{t,h}(c).
\end{equation}
两种写法条件层级不同。当前粒子实际保存抽取的 $e_t$，由粒子汇总不确定性；惯性存在时不能直接用 $\pi_t^-$ 替代 $\psi_t$。

示例不额外加入输出 lapse，也不对汇总类别概率再做置信锐化。作答变异来自知觉、执行和规则概率，这使口述来源保持清楚。

\section{口述模块：从本题判断到说出来的话}
\label{sec:oral}

首版 oral 在提交 \texttt{8019f29e} 中接入，配置名为 \texttt{rule\_content\_v1}。它允许省略、模糊命名和简单补充，是表达观测层，不另改搜索、信念或精度学习。

\subsection{先选择，再口述，再收到反馈}

口述读取 decision 保存的规则来源、类别概率和内部刺激快照。M1 一条路径有一条实际执行规则，因此选择与口述沿用同一 $e_t$，不另抽规则，不读取尚未收到的反馈或真实答案。

给定规则和已经选择的类别，先读取该类别区域的局部判断，不要求说出所有类别的完整方案。例如局部条件为“脖子长”和“脖子长于尾巴”，本题可能只说一句。

所选类别有分离区域时，读取距内部刺激最近的部分，并列最小则等权。即使选择不是规则最可能类别，也按实际选择的类别生成报告，这是明确的首版假设。

\subsection{步骤一：整理局部内容}

局部判断记为 $D_m=\{d_1,\ldots,d_K\}$，$m$ 为选中分量。带状区域的上下界合成一个带状判断，避免把同一接近关系误作两次独立报告。

词表由目录及允许的表达组合预先生成，不依赖真实口述、当前刺激或参数值。二分类390个代码，四分类1,498个，包括 \texttt{EMPTY/OTHER}；代码增多不表示规则增多。

\subsection{步骤二：每个核心判断可以省略}

每个判断独立地以 $q$ 被说出、以 $1-q$ 省略。两个判断都说、只说第一个、只说第二个、都不说的概率为
\begin{equation}
q^2,\quad q(1-q),\quad q(1-q),\quad(1-q)^2.
\end{equation}
$q=0.8$ 时为0.64、0.16、0.16、0.04。独立省略是模型假设，不表示每个人已被测到相同表达方式。

\texttt{EMPTY} 是最终没有表达内容。核心全省略后仍可能有补充句，所以不一定为空；记录缺失也不能直接当作 \texttt{EMPTY}。

\subsection{步骤三：核心单特征判断可以用躯干作模糊参照}

单特征、单方向阈值判断可改说“比躯干长／短”。该判断阈值为 $\theta_j$，命名概率为
\begin{equation}
b_j=\rho_B\exp[-(\theta_j-\mu_B)^2/(2s_B^2)],
\qquad\rho_B\in[0,1],\quad s_B>0.
\label{eq:body-wording}
\end{equation}
不说、普通表达、躯干表达的概率依次为
\begin{equation}
1-q,\qquad q(1-b_j),\qquad qb_j.
\end{equation}
不适用的判断令 $b_j=0$。比较方向保留，不把“比躯干长”改成短，也不把复杂关系任意改成单特征阈值。

$\mu_B,s_B$ 描述共同表达参照，$\rho_B$ 是命名倾向。物理躯干长度、行为阈值0.5与心理表达参照分开。躯干措辞不改变分类边界，当前也未拟合个人参照。

\subsection{步骤四：可以补充一个未参与局部规则的部位}

在非 \texttt{OTHER} 分支，以 $\alpha_O$ 尝试补充一句。特征须在完整局部判断中从未使用，且内部值不恰为0.5；在可用特征间等权选择一个，大于0.5说长，小于0.5说短。

“未使用”按完整规则算，不能因核心某句省略就把其特征当作补充候选。没有可用特征时不补充；最多补充一个普通长短句，不加任意复杂关系，补充也不改成躯干命名。

原规则使用脖子和尾巴，省略了“脖子长于尾巴”后再描述腿短，可能说“脖子长，而且腿短”。腿被提到却没有分类作用。相同句子也可能来自真正使用脖子和腿的规则，一次口述通常不能唯一确定来源。

\subsection{完整报告概率}

令 $\mathcal B(h,y,\widetilde{\bm x})$ 为并列最近的类别分量集合，$K_m(v)$ 为省略与命名后核心报告 $v$ 的概率，$U_m$ 为可补充特征，$a_j$ 为普通长短句，$\operatorname{code}$ 为规范组合代码。

$U_m$ 非空时，添加核为
\begin{equation}
H_m(o\mid v,\widetilde{\bm x})=
(1-\alpha_O)\ind[o=\operatorname{code}(v)]
+\frac{\alpha_O}{|U_m|}\sum_{j\in U_m}
\ind[o=\operatorname{code}(v\cup\{a_j\})].
\end{equation}
无可用特征时 $H_m=\ind[o=\operatorname{code}(v)]$。以 $\eta_O$ 生成预定义的 \texttt{OTHER}，最终
\begin{align}
\R(o\mid h,y,\widetilde{\bm x})
={}&\eta_O\ind(o=\texttt{OTHER})\nonumber\\
&+\frac{1-\eta_O}{|\mathcal B|}\sum_{m\in\mathcal B}\sum_v
K_m(v)H_m(o\mid v,\widetilde{\bm x}).
\label{eq:oral-kernel}
\end{align}
每步归一化。存在可用特征时，总补充事件概率为 $(1-\eta_O)\alpha_O$，示例0.02和0.2得到0.196。

\texttt{OTHER} 是预定义词表外内容汇总，不能自动接收所有无法解析、记录可疑的文本。未知代码不自动获得其概率，真实文字仍需编码和人工审查。

\subsection{选择与口述共享规则来源}

更一般地，令 $w_t(h)$ 为 decision 保留的实际来源权重：
\begin{equation}
P(y_t,o_t\mid\widetilde{\bm x}_t)=
\sum_h w_t(h)p_{t,h}(y_t)
\R(o_t\mid h,y_t,\widetilde{\bm x}_t).
\label{eq:choice-oral}
\end{equation}
给定实际选择，口述概率为联合质量除以 $\sum_h w_t(h)p_{t,h}(y_t)$。当前 M1 路径中 $w_t(h)=\ind(h=e_t)$，报告沿用执行规则。其他接口若使用规则混合，选择后贡献按 $w_t(h)p_{t,h}(y_t)$ 条件化，不把报告反写为学习信念。

先算“同一规则做出该选择并说出该报告”，再加总来源。分别平均选择与口述后相乘会丢失共享关系。

\Needspace{13\baselineskip}
\subsection{表达参数、接口与当前范围}

6个共同表达参数如下，例值均来自配置，尚未校准或按被试估计：

{\small
\begin{longtable}{L{1.6cm}L{8.0cm}L{2.1cm}L{1.5cm}}
\toprule
参数 & 含义 & 范围 & 例值\\
\midrule
\endhead
$q$ & 每个核心判断说出概率。 & $[0,1]$ & 0.8\\
$\alpha_O$ & 有可用特征时，非 \texttt{OTHER} 分支的补充概率。 & $[0,1]$ & 0.2\\
$\eta_O$ & 预定义词表外内容概率。 & $[0,1]$ & 0.02\\
$\rho_B$ & 核心单特征的躯干命名倾向。 & $[0,1]$ & 0.5\\
$\mu_B$ & 共同表达参照中心。 & $[0,1]$ & 0.75\\
$s_B$ & 共同表达参照宽度。 & $(0,\infty)$ & 0.1\\
\bottomrule
\end{longtable}
}

它们描述怎样表达，不描述怎样学习规则。$\rho_s$ 表示退出、$\rho_B$ 表示命名、$\alpha_O$ 表示补充，避免符号混用。

接口顺序为：
\begin{quote}
\small
\path{begin_trial}\ $\to$\ \path{predict_choice}\ $\to$\ 确定选择\\
$\to$\ \path{predict_oral}\ $\to$\ \path{complete_trial}。
\end{quote}
口述不更新学习、不消耗学习随机流；抽样用独立口述流。口述绑定选择后不再重抽 decision。

目前接通单条路径显式接口。行为 PF、批量拟合和自主生成不会自动读取口述，也没有用它修正粒子权重。输出 lapse、混合后置信变换和学习型标签映射需要另定义联合关系；当前口述接口拒绝这些不支持的组合。

\section{单试次顺序：哪些信息何时使用}
\label{sec:trial}

完整步骤如下：
\begin{enumerate}
\item 读取截至上一题的反馈与状态。第一题初始化；之后用上一题反馈递推 $F_t,E_t,g_t$。
\item 形成内部刺激 $\widetilde{\bm x}_t$；候选管理不使用本题真实类别或未来反馈。
\item 抽取非执行候选退出，条件化支持；按 $E_t$ 决定是否提出一个新人，完成竞争和支持迁移。
\item 幸存者保留精度，新人恢复初值；condition 3 同步设置联合配对先验。
\item 按信念与惯性抽取 $e_t$，用当前刺激和精度计算反馈前类别预测。
\item 自主生成抽取选择；拟合读取实际选择并修正粒子权重。需要口述时调用接口，使用同一 decision 快照。
\item 收到反馈后计算活跃规则证据；condition 3 计算联合证据。
\item 更新证据记忆，用缓存的反馈前证据更新精度，保存反馈供下一题搜索控制读取。
\end{enumerate}
失败递推在代码中于下一题开始时消费保存的反馈，与式\eqref{eq:failure}的概念时序一致。本题新反馈不能改变已经做出的本题预测或报告。反馈分别改变证据、精度和失败状态，不在相同状态里重复当成两条独立观测。

\section{粒子滤波：推断可能的内部路径}
\label{sec:pf}

\subsection{行为预测与路径权重}

设 $L$ 个粒子的作答前权重为 $w_{t,l}^-$，路径预测为 $p_{t,l}(c)$：
\begin{equation}
\widehat P_t(c)=\sum_{l=1}^{L}w_{t,l}^-p_{t,l}(c).
\label{eq:pf-prediction}
\end{equation}
实际选择后，
\begin{equation}
w_{t,l}^{\rm choice}\propto w_{t,l}^-p_{t,l}(y_t).
\end{equation}
能解释选择的路径获得高权重；随后各路径按实际选择与反馈学习。路径权重是研究者对路径的判断，路径内规则信念是被试对候选的判断，不能互换。

例如，两条路径原来各占一半，分别以0.8和0.2预测被试实际选择的答案，总体预测为0.5。观察这个选择后，两条路径的权重变为0.8和0.2。这表示第一条内部过程更能解释作答，不是把被试在每条路径内的规则信念都直接改成0.8和0.2。

有效样本量 $\mathrm{ESS}_t=1/\sum_l(w_{t,l}^{\rm choice})^2$。权重过度集中则按配置阈值重采样，保留候选、信念、精度、执行、失败与配对状态，再给未来随机流独立种子，不重新初始化学习。

\subsection{以后接入口述必须先汇总联合质量}

正确汇总为
\begin{equation}
\widehat P_t(y,o)=\sum_lw_{t,l}^-p_{t,l}(y)\R_l(o\mid y).
\label{eq:pf-oral-joint}
\end{equation}
给定选择的报告概率为它除以 $\widehat P_t(y)$，不能按选择前权重直接平均条件报告，因为实际选择会改变路径贡献。

式\eqref{eq:pf-oral-joint}是未来接入的关系，当前批量 PF 未执行。若用口述修正路径权重，需先明确编码、时点和缺失处理，再检查表达与学习参数的识别。

\subsection{怎样读懂汇总的候选和执行状态}

某规则出现在候选集合中的概率，可以按路径权重汇总其是否活跃：
\begin{equation}
\widehat P_t(h\in A_t)=\sum_lw_{t,l}^-\ind(h\in A_{t,l}),
\qquad
\overline N_t=\sum_lw_{t,l}^-|A_{t,l}|.
\end{equation}
每条路径的候选数量是整数，汇总平均却可以是2.3。这表示不同可能路径的候选数量不同，不表示被试实际保留了2.3条规则。某规则的执行概率也按各路径是否执行它汇总，与候选可访问概率和规则信念分别解释。

使用反馈前权重得到的是本题预测时可用的状态；观察选择后再汇总，得到的是已结合本题选择的在线推断。分析图应明确是哪一种，不能混用时点来说明某个状态提前导致了行为变化。

\Needspace{30\baselineskip}
\section{个体参数估计：估计什么，怎样评价}
\label{sec:parameters}

\subsection{行为核心的14个候选被试参数}

搜索、候选与执行共10项，另加记忆1项、精度3项：
\begin{align}
\theta_s=(&E_{0,s},a_s,\delta_s,g_{0,s},c_{G,s},
M_s^{\max},\rho_s,\omega_s,\tau_{T,s},\lambda_s,\nonumber\\
&\gamma_s,\beta_{0,s},\eta_{+,s},\eta_{-,s}).
\end{align}
这是待估计清单，不代表14项已能独立识别，也不要求每次同时自由估计全部。自由项、范围和固定项须拟合前规定。下表例值来自三条件 M1 配置，只用于说明当前可执行设置，各字段可按被试覆盖。

{\small
\begin{longtable}{L{1.8cm}L{6.8cm}L{3.4cm}L{1.5cm}}
\toprule
参数 & 认知含义 & 数学范围 & 例值\\
\midrule
\endhead
$E_{0,s}$ & 基础搜索概率。 & $(0,1)$ & 0.05\\
$a_s$ & 失败推动搜索的增益。 & $[0,\infty)$ & 1.0\\
$\delta_s$ & 失败状态保持。 & $[0,1)$ & 0.6\\
$g_{0,s}$ & 基础全局比例。 & $[0,1]$ & 0.2\\
$c_{G,s}$ & 失败对范围的增益。 & $[0,1]$ & 0.5\\
$M_s^{\max}$ & 候选数量上限。 & 整数；\newline $2\leq M_s^{\max}<J$ & 3\\
$\rho_s$ & 备选退出率。 & $[0,1]$ & 0.1\\
$\omega_s$ & 局部结构来源倾向。 & $[0,1]$ & 0.5\\
$\tau_{T,s}$ & 功能支持泛化尺度。 & $(0,\infty)$ & 0.1\\
$\lambda_s$ & 上一执行规则的惯性。 & $[0,\infty)$ & 1.0\\
$\gamma_s$ & 规则及联合配对证据保持。 & $[0,1]$ & 0.8\\
$\beta_{0,s}$ & 新规则初始精度。 & $[\beta_{\min},\beta_{\max}]$ & 5.0\\
$\eta_{+,s}$ & 精度增长速率。 & $[0,1]$ & 0.04\\
$\eta_{-,s}$ & 精度下降速率。 & $[0,1]$ & 0.15\\
\bottomrule
\end{longtable}
}
知觉参数来自独立测量；6个表达参数目前单列为共同演示设置，不加入已经实施的行为拟合。配对权重、信念、失败、实际数量和执行规则均为状态。

\subsection{用反馈前选择概率评价参数}

行为目标可用平均选择负对数概率：
\begin{equation}
\mathrm{NLL}_{\rm choice}(\theta_s)=
-\frac1{T_{\rm score}}\sum_{t\in\mathcal T_{\rm score}}
\log\widehat P_t(y_t;\theta_s).
\label{eq:nll}
\end{equation}
保存计分集合、顺序、训练／评价边界与随机预算，不能拿反馈后预测替代反馈前概率。多次运行应先逐题平均类别概率再评分；平均各次 NLL 是另一统计量。

粒子与重复表达隐藏路径或数值不确定性，不自动给出参数置信区间。完整行为序列选出的参数在同序列上仍是样本内拟合，即使状态递推没有使用未来反馈。

\subsection{正式搜索与识别尚需确定}

现有工具支持完整配置路径覆盖和候选比较。拟合可按记忆、搜索、候选、精度等相关块组织，共同随机数比较后再用独立种子检验稳定性。M1 正式范围、网格、粒子预算、停止与恢复协议尚未冻结，不把 smoke 配置称为正式方案。

失败增益小时 $\delta_s$ 难识别，$g_0$ 接近1时 $c_G$ 作用弱，两局部核相近或大量回退时 $\omega_s$ 难识别，少到达上限时容量可能只得宽范围。支持泛化、惯性、精度与记忆也可能共同解释行为，需恢复和针对性比较。

达到与0826相当效果，要在可比数据、计分试次和预算下分别比较选择、学习过程与口述预测。测试和短序列不能证明表现保持。

\section{Condition 3：层级反馈与配对学习}
\label{sec:pairing}

\subsection{知道层级，不知道具体反应配对}

被试知道两个科、每科两个种，但不知道哪两个反应同科。三种配对为
\begin{equation}
G_1=\{\{1,2\},\{3,4\}\},\quad
G_2=\{\{1,3\},\{2,4\}\},\quad
G_3=\{\{1,4\},\{2,3\}\}.
\end{equation}
开始等权。数据 \texttt{category} 为真实类别、\texttt{choice} 为统一反应编号、\texttt{presskey} 为实际按键；转换表用于读懂数据，不预告正确配对。

\subsection{同次反馈评价规则和配对}

每规则有 $q_t^-(k\mid h)$，联合信念为
\begin{equation}
\Pi_t^-(h,k)=\pi_t^-(h)q_t^-(k\mid h),\qquad q_1^-(k\mid h)=1/3.
\end{equation}
令 $s_k(y)$ 为配对中同组的另一个反应：
\begin{equation}
\ell_t(h,k)=
\begin{cases}
p_{t,h}(y_t),&r_t=1,\\
p_{t,h}[s_k(y_t)],&r_t=0.5,\\
1-p_{t,h}(y_t)-p_{t,h}[s_k(y_t)],&r_t=0.
\end{cases}
\label{eq:ternary-feedback}
\end{equation}
完全正确支持所选反应，半分支持同科另一种，0分支持另科两个反应。事件互斥且概率和为1，0分不能用二值错误概率。

联合更新为
\begin{equation}
\Pi_t^+(h,k)=
\frac{[\Pi_t^-(h,k)]^{\gamma_s}\ell_t(h,k)}
{\sum_{v\in A_t}\sum_{j=1}^3[\Pi_t^-(v,j)]^{\gamma_s}\ell_t(v,j)}.
\label{eq:joint-memory}
\end{equation}
规则边际为 $\sum_k\Pi_t^+(h,k)$，条件配对在规则内归一化。复用 $\gamma_s$，不增加配对遗忘率；衰减联合信念不同于分别衰减两套边际。

配对均匀时，半分和0分的规则证据分别为 $[1-p_{t,h}(y_t)]/3$ 和 $2[1-p_{t,h}(y_t)]/3$，相对规则支持相同，配对支持却可以不同。配对学会后相同反馈给出更具体线索，不人为规定第几题开始使用层级。

\subsection{三值反馈怎样更新精度}
\label{sec:pairing-precision}

遗忘后、加入本题反馈前的条件配对为
\begin{equation}
\widetilde q_t^-(k\mid h)=
\frac{[\Pi_t^-(h,k)]^{\gamma_s}}{\sum_j[\Pi_t^-(h,j)]^{\gamma_s}}.
\end{equation}
活跃联合单元在 $\gamma_s=0$ 时按等权处理；某规则行的质量为0时，其条件配对使用均匀分布，避免未定义的除法。
边际反馈概率 $b_t(h)=\sum_k\widetilde q_t^-(k\mid h)\ell_t(h,k)$。均匀四类下，正确、半分、0分的机会概率为 $1/4,1/4,1/2$，支持量为
\begin{equation}
\kappa_t^{(3)}(h)=
\frac{b_t(h)-b_{\rm chance}(r_t)}
{b_t(h)+b_{\rm chance}(r_t)}.
\label{eq:ternary-precision-support}
\end{equation}
进入式\eqref{eq:precision-update}。均匀预测三种反馈都中性，半分可以支持或反驳。代码缓存反馈前绝对证据，避免记忆更新后再次用新配对评价同一反馈。搜索则仍用 $1-r_t$，服务不同认知问题。

\subsection{候选变化后保留配对知识}
\label{sec:pairing-transfer}

迁移前保存整体边际 $\mu_t(k)=\sum_{h\in A_t}\Pi_t^+(h,k)$。幸存者保留条件配对，新人继承 $\mu_t$：
\begin{equation}
\Pi_{t+1}^-(h,k)=\pi_{t+1}^-(h)
\begin{cases}
q_t^+(k\mid h),&h\text{幸存},\\
\mu_t(k),&h\text{为新人}.
\end{cases}
\end{equation}
改变分类规则不必从头猜反应关系，但新人仍不知道真实配对；局部／全局来源不改变此约定。规则权重变化后，整体配对边际可能变化，不能称全部联合状态守恒。

配对集中不等于配对正确，更不等于会按刺激分科。科、种、平均分分别评价；口述如何表达配对知识目前没有独立机制。

\section{从模块状态到认知机理解释}
\label{sec:interpretation}

\subsection{发现、接纳和执行分开看}

相近正确率可能来自不同过程：未想到有效规则，想到后未入场，入场后支持低，支持提高后未执行，或已执行但精度低。分析时同时看可访问性、信念、执行和精度，不能只看最高概率规则。

这些是拟合条件下推断的隐藏过程，需要行为和口述约束，不因模型设了变量就表示已识别。

\subsection{候选数量变化的含义}

连续成功使失败回落、搜索减少，弱备选仍可退出，集合可能收缩到一条。再次失败可以增加新人，集合扩张。低支持未退出者仍能学习和执行，已退出者需重新想到；“接近0”和“不再考虑”不是同一状态。

\subsection{结构与功能变化分别量化}

只改一个特征可改变很多分类；分类变化小也不保证只改一处判断。分析编辑类型、功能距离和来源，有助于区分调整思路与想到功能相近的另一解释。共同距离与成本是可复现定义，还不是个人心理测量。

\subsection{口述怎样约束解释}

口述可检查内容是否变化、判断是否反复出现、表达变化与执行是否同步。省略和补充使同一句有多种来源，不能直接确定唯一规则或把新提到的部位当作规则扩张。

后续机理图可以围绕候选进入、支持积累、执行切换、精度增长与表达变化的时间关系组织，并呈现粒子不确定性。快速、渐进、晚期突变是待解释现象，不预设为被试身份输入。本节提出分析方向，没有宣称已经重现这些现象。

\section{共同设置、代码对应与验证边界}
\label{sec:implementation}

\subsection{固定定义与试次状态}

共同定义包括目录与标签、均匀先验、初始一条、单新人、执行保护、退出函数、最低临时权重竞争、编辑语法与成本、均匀参考分布、100,000点与种子0、$\tau_P=0.10$、精度上下界、反馈增益1、永久记忆权重0。口述6项目前也是共同演示设置。

固定项仍是模型假设。改变参考分布、成本、标签或初始化须记录版本与比较，不悄悄修正某些被试。实际数量、集合、信念、精度、失败、执行与配对为递推状态；粒子、重复和种子为数值预算。

\subsection{主要代码和配置}

机制位于共享 \path{src/Bayesian_state/}，下面文件相对该目录：

{\small
\begin{longtable}{L{6.6cm}L{7.0cm}}
\toprule
位置 & 对应内容\\
\midrule
\endhead
\path{model/modules/perception.py} & 知觉。\\
\path{hypothesis_space/geometry/boundary.py} & 类别区域距离。\\
\path{model/modules/memory.py} & 二值反馈证据记忆。\\
\path{model/modules/pairing_memory.py} & 配对学习与迁移。\\
\path{model/modules/hypothesis_transition/}\newline
\path{mechanistic_rule_search.py} & 搜索控制、候选与执行。\\
\path{hypothesis_space/rule_programs.py} & 判断程序和合法编辑。\\
\path{hypothesis_space/similarity.py} & 类别对应张量与一致率。\\
\path{model/modules/beta.py} & 精度学习。\\
\path{model/readout.py} & 选择与来源快照。\\
\path{model/oral_readout.py}\newline
\path{model/oral_report_space.py}\newline
\path{model/oral_observation.py} & 口述表达、词表与联合概率。\\
\path{model/state_model.py} & 试次接口和时序。\\
\bottomrule
\end{longtable}
}

三条件配置位于 \path{configs/exp123/model_struct/}：
\begin{itemize}
\item \path{model_0923_cond1_M1.yaml}；
\item \path{model_0923_cond2_M1.yaml}；
\item \path{model_0923_cond3_M1.yaml}。
\end{itemize}
搜索／候选位于 \path{modules.hypo_transitions_mod.kwargs}，搜索五项在其 \path{search_controller} 下，记忆、精度、口述分别在自身配置块。

科学代码基于 \texttt{4c1ccd79}，oral 来源为 \texttt{8019f29e}。配套 \path{model_0923.md} 保留讨论与历史，本文按 M1 连贯说明。

\subsection{已检查什么，尚未证明什么}

2026年10月10日297项相关测试通过，覆盖程序与几何、两局部核和回退、张量、失败保持、数量收缩、容量与质量守恒、精度重置、配对继承、执行／口述共用、随机重现、粒子重采样和时序。旧回归包含期刊独立数值参考，没有重新生成参考。

三个条件分别用 S122、S206、S314，各32题、8粒子、一次重复、一个 worker 做真实数据流程检查，通过概率归一化、容量、反馈时序与配对分布检查。输出在独立临时目录。环境为 Python 3.11.10、NumPy 1.26.4、SciPy 1.13.1、PyYAML 6.0.2。

这说明实现和短流程可运行，不证明例值拟合得好。完整拟合、恢复、口述校准、真实联合评分和与0826的可比效果检查仍未完成；成本、退出、初始化、惯性和支持泛化的心理意义仍需审视。

\subsection{本文的生成与复核}

源文件为 \path{model_0923.tex}，可读文件为 \path{model_0923.pdf}。运行 \path{compile_model_0923.sh} 用 XeLaTeX 或 XeTeX 的 LaTeX 格式编译三遍，使目录和交叉引用稳定。日志保留在新的系统临时目录；脚本不覆盖已有 PDF，重编译时指定新输出路径。

本文与实际配置一起阅读。例值帮助理解和检查，不是被试估计或正式冻结范围。

\clearpage
\appendix
\section{四分类目录的完整分组}
\label{app:catalogue}

下表按实际家族计数。交叉表示同时使用两个判断，条件分支表示先判断、再在相应分支区分。

{\small
\begin{longtable}{L{5.5cm}L{6.7cm}L{1.3cm}}
\toprule
代码家族 & 内容 & 数量\\
\midrule
\endhead
\texttt{2d\_axis\_pair} & 两个单特征阈值交叉。 & 6\\
\texttt{2d\_equality\_sum} & 同一特征对的顺序与和交叉。 & 6\\
\texttt{3d\_axis\_equality} & 单特征阈值与另两特征顺序。 & 12\\
\texttt{3d\_axis\_sum} & 单特征阈值与另两特征和。 & 12\\
\texttt{3d\_axis\_triple} & 三阈值的条件分支。 & 24\\
\texttt{3d\_axis\_equality\_sum} & 阈值分支下顺序或和。 & 24\\
\texttt{4d\_equality\_pair} & 不重叠特征对的顺序交叉。 & 3\\
\texttt{4d\_sum\_pair} & 不重叠特征对的和交叉。 & 3\\
\texttt{4d\_equality\_axis\_pair} & 顺序分支与剩余特征阈值。 & 12\\
\texttt{4d\_sum\_axis\_pair} & 和分支与剩余特征阈值。 & 12\\
\texttt{dimension\_max/min} & 最大／最小特征分类。 & 2\\
\texttt{axis\_pair\_overlap} & 阈值与共享一特征的顺序交叉。 & 12\\
\midrule
合计 & & 128\\
\bottomrule
\end{longtable}
}
程序用于表示审计与编辑邻居，实际类别预测调用共享几何。程序与几何的标签对应已核对，不表示目录已覆盖所有真实口述。
\end{document}
